\section{Derivation of the General Emission Integral}
\label{app:general_emission}

Consider an electronic transition between an initial state $\ket{\mathbf{k}_{1}, \alpha}$ and a final state $\ket{\mathbf{k}_{2}, \beta}$, mediated by an energy transfer $\hbar \omega$. The subscripts $\alpha$ and $\beta$ represent specific energy bands (valence or conduction), allowing transitions between different branches of the dispersion relation. According to Fermi's Golden Rule, the microscopic transition rate is
\begin{equation}
    \Gamma_{1\to 2}(\hbar\omega) = \frac{2\pi}{\hbar} \left|\mu_{12}^{\alpha\beta}\right|^{2} \ \delta\left(\E_{\beta}(\mathbf{k}_{2}) - \E_{\alpha}(\mathbf{k}_{1}) + \hbar\omega \right)
\end{equation}
where $\mu_{12}^{\alpha\beta}$ is the transition dipole moment governing momentum conservation, and the Dirac delta function enforces energy conservation based on the band-specific dispersion relations $\E_{\alpha}$ and $\E_{\beta}$.

For clarity, we temporarily suppress the band indices to derive the rate for a generic pair of bands. To determine the rate $\Gamma_{1}$ from a specific initial state $\ket{\mathbf{k}_{1}}$, we sum over all available final states $\ket{\mathbf{k}_{2}}$. This summation is weighted by the probability of the final state being vacant, $[1-f(\mathbf{k}_{2})]$, and accounts for spin degeneracy. In the continuum limit, the discrete summation over wave-vectors transforms into an integral over the Brillouin Zone (BZ) with a density of states factor $V/(2\pi)^{3}$:
\begin{align}
    \Gamma_{1}(\hbar\omega) &= 2 \cdot \sum_{\mathbf{k}_{2}} \ \Gamma_{1\to 2}(\hbar\omega) \Big[1-f(\mathbf{k}_{2})\Big] \\
    &= 2V \int_{\text{BZ}} \frac{d^{3}k_{2}}{(2\pi)^{3}} \ \Gamma_{1\to 2}(\hbar\omega) \Big[1-f(\mathbf{k}_{2})\Big]
\end{align}
Similarly, to obtain the total rate for this band pair, we sum the single-state rates over all possible initial states $\ket{\mathbf{k}_{1}}$, weighted by the occupation probability $f(\mathbf{k}_{1})$
\begin{align}
    \Gamma (\hbar\omega) &= 2\cdot \sum_{\mathbf{k}_{1}} \ \Gamma_{1}(\hbar\omega) f(\mathbf{k}_{1}) \\
    &= 2V \int_{\text{BZ}} \frac{d^{3}k_{1}}{(2\pi)^{3}} \ \Gamma_{1}(\hbar\omega) \ f(\mathbf{k}_{1})
\end{align}
Finally, to account for all possible emission channels, we reintroduce the band indices and perform a general summation over all bands $\alpha$ and $\beta$. This covers all permutations, including interband (conduction-valence $\alpha\beta = cv$) and intraband (conduction-conduction $\alpha\beta =cc$) transitions. Defining $\overline{f_{\beta}(\mathbf{k})} \equiv 1 - f_{\beta}(\mathbf{k})$, the general expression is
\begin{equation}
    \Gamma(\hbar\omega) = \frac{2\pi}{\hbar} \left( \frac{2V}{(2\pi)^{3}} \right)^{2} \sum_{\alpha, \beta} \iint_{\text{BZ}} |\mu^{\alpha\beta}(\mathbf{k_{1}},\mathbf{k_{2}})|^{2} \ f_{\alpha}(\mathbf{k}_{1})\overline{f_{\beta}(\mathbf{k}_{2})} \ \delta\big(\E_{\beta}(\mathbf{k}_{2}) - \E_{\alpha}(\mathbf{k}_{1}) + \hbar\omega \big) \ d^{3}k_{1} \ d^{3}k_{2} \label{eq:general_integral}
\end{equation}

\section{Derivation of Analytic Approximation}
\label{app:analytic}

Starting from the general intraband rate ($\alpha=\beta=c$) with constant matrix element approximation $|\mu|^{2} \to |\mu_{cc}|^{2}$:
\begin{equation}
    \Gamma_{\text{cc}}(\hbar\omega) = \frac{2\pi}{\hbar}\left( \frac{2\mu_{\text{cc}}V}{(2\pi)^{3}} \right)^{2} \iint_{\text{BZ}} f\Big[\E_{\text{c}}(\mathbf{k}_{\text{i}})  \Big]\overline{f\Big[\E_{\text{c}}(\mathbf{k}_{\text{f}})  \Big]} \ \delta \Big( \E_{\text{c}}(\mathbf{k}_{\text{i}}) - \E_{\text{c}}(\mathbf{k}_{\text{f}}) - \hbar\omega \Big) \ d^{3}k_{\text{i}}d^{3}k_{\text{f}}
\end{equation}
Assuming isotropic parabolic bands allows transforming from $d^{3}k$ to $\rho(\E)d\E$:
\begin{equation}
    \Gamma_{\text{cc}}(\hbar\omega) = C \iint f(\E_{i})\overline{f(\E_{f})} \delta(\E_{i} - \E_{f} - \hbar\omega) \rho(\E_{i})\rho(\E_{f}) d\E_{i} d\E_{f}
\end{equation}
where $C$ collects the prefactors. Resolving the delta function $\E_{f} = \E_{i} - \hbar\omega$:
\begin{equation}
    \Gamma_{\text{cc}}(\hbar\omega) = C \int_{0}^{\infty} f(\E)\left[1 - f(\E-\hbar\omega)\right] \rho(\E)\rho(\E-\hbar\omega) d\E
\end{equation}
Changing variables $\E \to \E + \hbar\omega$:
\begin{equation}
    \Gamma_{\text{cc}}(\hbar\omega) = C \int_{0}^{\infty} f(\E+\hbar\omega)\left[1 - f(\E)\right] \rho(\E+\hbar\omega)\rho(\E) d\E
\end{equation}
Assuming constant density of states $\rho(\E) \approx \rho(\E_{F})$ near the Fermi level:
\begin{equation}
    \Gamma_{\text{cc}}(\hbar\omega) \approx C \rho^{2}(\E_{F}) \int_{0}^{\infty} f(\E+\hbar\omega)\left[1 - f(\E)\right] d\E
\end{equation}
Using the identity $\feq(\E+\hbar\omega)[1-\feq(\E)] = n_{B}(\hbar\omega) [\feq(\E) - \feq(\E+\hbar\omega)]$, the integral can be evaluated analytically for equilibrium distributions:
\begin{equation}
    \int_{0}^{\infty} \feq(\E+\hbar\omega)[1-\feq(\E)] d\E \approx \frac{\hbar\omega}{e^{\beta\hbar\omega}-1}
\end{equation}
yielding the final analytic approximation:
\begin{equation}
    \Gamma_{\text{cc}}(\hbar\omega) \propto \frac{\hbar\omega}{e^{\hbar\omega/k_{B}T}-1}
\end{equation}

\section{Interband Transitions and Joint Density of States}
\label{app:interband_jdos}

For interband transitions, momentum conservation ($\mathbf{k}_{c} = \mathbf{k}_{v} \equiv \mathbf{k}$) simplifies the double integral. The transition dipole becomes diagonal $\delta(\mathbf{k}_{1}-\mathbf{k}_{2})$:
\begin{equation}
    \Gamma(\hbar\omega) \propto \int_{\text{BZ}} |\mu^{cv}(\mathbf{k})|^{2} f_{c}(\mathbf{k})\overline{f_{v}(\mathbf{k})} \delta(\E_{c}(\mathbf{k}) - \E_{v}(\mathbf{k}) - \hbar\omega) d^{3}k
\end{equation}
Defining $\E_{cv}(\mathbf{k}) = \E_{c}(\mathbf{k}) - \E_{v}(\mathbf{k})$ and introducing the Joint Density of States (JDOS):
\begin{equation}
    \rho_{J}(\E) = \frac{2}{(2\pi)^{3}} \int_{\text{BZ}} \delta(\E_{cv}(\mathbf{k}) - \E) d^{3}k
\end{equation}
For parabolic bands with $\E_{c} = \E_{g} + \hbar^{2}k^{2}/2m_{c}$ and $\E_{v} = -\hbar^{2}k^{2}/2m_{v}$, using reduced mass $\mu_{r}^{-1} = m_{c}^{-1} + m_{v}^{-1}$:
\begin{equation}
    \rho_{J}(\E) = \frac{1}{2\pi^2} \left( \frac{2\mu_r}{\hbar^2} \right)^{3/2} \sqrt{\E - \E_g}
\end{equation}

\section{Electronic Density of States}
\label{app:edos}
For a free electron gas in 3D, the transformation from k-space to energy space is given by:
\begin{equation}
    \rho(\E) = \frac{V}{2\pi^{2}} \left(\frac{2m}{\hbar^{2}}\right)^{3/2} \sqrt{\E}
\end{equation}
This derivation follows from $\E = \hbar^{2}k^{2}/2m \implies k = \sqrt{2m\E}/\hbar$ and $d^{3}k = 4\pi k^{2} dk$.

\section{Thermal Factor in Energy Space}
\label{app:thermal_factor}
The thermal factor $I = \int f(\E+\hbar\omega)[1-f(\E)] d\E$ governs the emission spectrum shape.
In equilibrium:
\begin{align}
    I &= n_{B}(\hbar\omega) \int_{0}^{\hbar\omega} f(\E) d\E \\
    &= n_{B}(\hbar\omega) \left[ \hbar\omega + \frac{1}{\beta}\ln(1+e^{-\beta\mu}) - \frac{1}{\beta}\ln(1+e^{\beta(\hbar\omega-\mu)}) \right]
\end{align}
In the limit $\mu \gg k_{B}T, \hbar\omega$, this simplifies to $\hbar\omega / (e^{\beta\hbar\omega}-1)$.

\section{Numeric Delta Approximation}
\label{app:numeric_delta}
To numerically evaluate integrals containing $\delta(x)$, we approximate it with a Gaussian $G_{\sigma}(x)$:
\begin{equation}
    \delta(x) \approx G_{\sigma}(x) = \frac{1}{\sigma\sqrt{2\pi}} \exp\left(-\frac{x^{2}}{2\sigma^{2}}\right)
\end{equation}
The approximation requires strict convergence checks:
\begin{enumerate}
    \item $\sigma \to 0$: The Gaussian width must be small enough to resolve features in the integrand.
    \item $N \to \infty$: The sampling density must be high enough to capture the Gaussian peak ($\Delta x \ll \sigma$).
    \item Integration limits: Must cover $\sim \pm k\sigma$ (e.g., $k=5$) to capture the Gaussian tail.
\end{enumerate}
We use a variable integration range that scales with $\sigma$, ensuring $N \propto 1/\sigma$ to maintain constant sampling quality.
